\chapter{Modello e stack software}
\label{cap:modello}

\section{Il modello base}

Il modello impiegato è \texttt{Qwen2.5-0.5B-Instruct}: un Transformer
decoder-only da circa $0{,}5$ miliardi di parametri, nella variante già
sottoposta a istruzione (\emph{instruct}), cioè addestrata a seguire richieste
formulate in linguaggio naturale.

La scelta di un modello piccolo è motivata da tre ragioni, in ordine di peso:

\begin{enumerate}
  \item \textbf{Vincolo di risorse.} L'addestramento gira su una singola GPU di
        un cluster condiviso con una quota di una sola allocazione per volta
        (Capitolo~\ref{cap:infrastruttura}). Il reinforcement learning richiede
        di generare $G$ candidati per ogni prompt a ogni passo, quindi il costo
        di generazione domina; un modello grande renderebbe l'esperimento
        infattibile.
  \item \textbf{Adeguatezza al compito.} Come stabilito nel
        Capitolo~\ref{cap:dataset}, il compito ha meno di un bit di incertezza
        residua per token. Non è un compito che richiede capacità di
        ragionamento: richiede di memorizzare una mappa lessicale.
  \item \textbf{Iterazione.} Un modello piccolo permette cicli di esperimento
        brevi, che è ciò che serve quando l'oggetto dell'indagine è il
        \emph{disegno} sperimentale e non la prestazione assoluta.
\end{enumerate}

La variante \emph{instruct} è preferita alla variante base perché il prompt
adottato ha forma di istruzione, e perché nella cella \emph{few-shot} gli esempi
sono inseriti come turni di conversazione.

\section{Adattamento: LoRA e quantizzazione}

Il modello base non viene aggiornato. Si applica l'adattamento a rango basso
descritto nella~\eqref{eq:lora} con la configurazione seguente.

\begin{center}
\begin{tabular}{lr}
\toprule
Iperparametro & Valore \\
\midrule
Rango $r$ & $32$ \\
Scala $\alpha$ & $64$ \\
Dropout & $0{,}05$ \\
Moduli adattati & $7$ \\
Quantizzazione dei pesi base & 4 bit (QLoRA) \\
\bottomrule
\end{tabular}
\end{center}

I sette moduli sono \texttt{q\_proj}, \texttt{k\_proj}, \texttt{v\_proj},
\texttt{o\_proj} (le quattro proiezioni dell'attenzione) e \texttt{gate\_proj},
\texttt{up\_proj}, \texttt{down\_proj} (le tre proiezioni del blocco
feed-forward). Adattare anche il blocco feed-forward, e non solo l'attenzione, è
la configurazione che la letteratura su QLoRA~\cite{dettmers2023qlora} indica
come preferibile quando il numero di parametri addestrabili non è il vincolo
stringente.

Il rapporto $\alpha/r = 2$ è la scelta convenzionale. Il dropout $0{,}05$ agisce
solo sul ramo dell'adattatore.

Il caricamento è realizzato in \texttt{src/models/model\_loader.py}, che
incapsula sia la quantizzazione a 4 bit sia l'installazione degli adattatori,
esponendo una sola funzione di caricamento usata da tutti i punti di ingresso
(SFT, GRPO, valutazione). Questa centralizzazione è deliberata: garantisce che
la cella di valutazione carichi il modello esattamente come la cella di
addestramento, evitando una classe di discrepanze difficile da diagnosticare.

\section{Lo stack software}

La riproducibilità richiede di fissare le versioni. Lo stack è il seguente.

\begin{center}
\begin{tabular}{lll}
\toprule
Componente & Versione & Ruolo \\
\midrule
PyTorch & \texttt{2.7.1+cu118} & tensori, autograd \\
CUDA & $11{,}8$ & runtime GPU \\
Transformers & $5{,}3{,}0$ & modello, tokenizzatore, generazione \\
TRL & $0{,}24{,}0$ & \texttt{SFTTrainer}, \texttt{GRPOTrainer} \\
PEFT & $0{,}19{,}1$ & implementazione LoRA \\
Unsloth & --- & kernel ottimizzati, caricamento 4 bit \\
torchao & $0{,}17{,}0$ & primitive di quantizzazione \\
sacrebleu & $2{,}6{,}0$ & BLEU e chrF \\
sentence-transformers & $5{,}2{,}3$ & retrieval per il few-shot \\
\bottomrule
\end{tabular}
\end{center}

Alcune di queste versioni non sono neutrali, e vale la pena segnalarlo qui
perché il Capitolo~\ref{cap:infrastruttura} documenta i guasti che ne sono
derivati.

\begin{description}
  \item[Unsloth.] Fornisce kernel ottimizzati e un percorso di caricamento a 4
        bit. Il suo import è però costoso e con effetti collaterali globali: va
        eseguito \emph{dopo} le validazioni di configurazione, non prima. Un bug
        reale documentato nella sezione~\ref{sec:bug-evalonly} consisteva
        esattamente in una validazione posta dopo l'import, con il risultato che
        un errore di configurazione banale si manifestava solo dopo minuti di
        caricamento.
  \item[Transformers 5.3.0.] Ha introdotto un cambiamento di firma in una
        funzione interna di rilevamento dei pacchetti, che ha rotto l'import di
        \texttt{trl.trainer.grpo\_trainer} (sezione~\ref{sec:bug-weave}).
  \item[TRL 0.24.0.] I valori predefiniti del \texttt{GRPOTrainer} in questa
        versione determinano quale variante dell'obiettivo GRPO viene
        effettivamente ottimizzata. Poiché nessuna configurazione storica del
        progetto li impostava esplicitamente, gli esperimenti hanno usato i
        default della libreria per omissione e non per scelta. L'analisi
        completa è nella sezione~\ref{sec:denominatori}.
\end{description}

\section{Il formato del prompt}

Il modello riceve un'istruzione e produce il gloss. Sono state usate due
modalità.

\subsection{Zero-shot}

Il prompt contiene solo l'istruzione e la frase da tradurre. La lunghezza
massima del prompt è fissata a $256$ token.

\subsection{Few-shot con retrieval}

Il prompt contiene, prima della frase da tradurre, $k=3$ esempi
\emph{recuperati} dal training set in base alla somiglianza semantica con la
frase corrente. Il recupero è realizzato con \texttt{sentence-transformers}: le
frasi di training sono codificate in vettori una volta sola, e a inferenza si
cercano i tre vicini più prossimi.

La lunghezza massima del prompt in questa modalità è $768$ token, per far posto
agli esempi. Questo vincolo è verificato automaticamente:
\texttt{tests/validate\_configs.py} rifiuta qualunque configurazione che attivi
il retrieval con \texttt{max\_prompt\_length} inferiore a $512$, perché un
prompt troncato eliminerebbe silenziosamente gli esempi rendendo la cella
indistinguibile dallo zero-shot.

\subsection{Perché questa distinzione è centrale}

Anticipando il Capitolo~\ref{cap:risultati}: la differenza fra zero-shot e
few-shot è, numericamente, il fattore \emph{più grande} osservato in tutto
l'esperimento. Il modello base con vincolo di decodifica passa da ROUGE-L
$0{,}1373$ (zero-shot) a $0{,}4664$ (few-shot), un salto di $0{,}33$ ottenuto
\emph{senza aggiornare un solo parametro}. Il totale di variazione fra il punto
di partenza peggiore e il miglior risultato con reinforcement learning è circa
$0{,}42$. Ne segue che la maggior parte dell'effetto attribuibile alla
``pipeline'' è in realtà effetto di prompting.

Questa è una confusione di fattori nel disegno originale: le celle che
combinavano metodo di addestramento e modalità di prompting non permettevano di
separare i due contributi. La correzione proposta nel
Capitolo~\ref{cap:conclusioni} è di rendere l'esperimento primario una griglia
prompting $\times$ decodifica, con il metodo di addestramento come fattore
successivo.

\section{Il vocabolario di uscita}

Il modello genera token, ma l'uscita valida è una sequenza di gloss. Il
collegamento fra i due livelli è il vocabolario gloss di $15\,472$ elementi
descritto nel Capitolo~\ref{cap:dataset}, tradotto in un insieme di $3\,967$
identificatori di token ammessi come \emph{primo} token di un gloss.

La costruzione, la doppia radice necessaria per gestire il token con spazio
iniziale, e le quattro forme gloss che risultano non emettibili sono trattate
nel Capitolo~\ref{cap:decoding}.
